% !TeX root = ../main.tex
\chapter*{Notation and Units}
\addcontentsline{toc}{chapter}{Notation and Units}
\markboth{Notation and Units}{Notation and Units}
\label{ch:notation}

\noindent\textbf{Sign convention.}\enspace Compression is positive. This is
standard in rock mechanics and it is the opposite of the structural
engineering convention, in which tension is positive. The difference matters
most at interfaces between disciplines --- when lining thrusts and moments
computed by a geotechnical model are passed to a structural check, or when a
finite element package written for structures is applied to ground. Where a
result is quoted from a structural source, the convention is stated
explicitly.

\noindent\textbf{Principal stresses.}\enspace Principal stresses are ordered
$\sigma_1 \geq \sigma_2 \geq \sigma_3$, so that $\sigma_1$ is the major and
$\sigma_3$ the minor principal stress. In failure criteria the pair
$(\sigma_1, \sigma_3)$ refers to the stresses acting on the element under
consideration, not to the in-situ field. In-situ stresses are written
$\sigma_v$ and $\sigma_h$, and the ratio $k = \sigma_h / \sigma_v$. Confusing
the in-situ field with the criterion stresses is among the most common errors
in parameter derivation, and Chapter~\ref{ch:rock-mechanics} treats it
explicitly.

\noindent\textbf{Effective stress.}\enspace A prime denotes an effective
stress: $\sigma' = \sigma - u$, where $u$ is pore pressure. Rock mass
parameters derived from the Hoek--Brown criterion are total-stress quantities
unless stated otherwise, and combining them with effective-stress soil
parameters in the same model requires care.

\noindent\textbf{Units.}\enspace SI units are used throughout. Stresses and
strengths are in \si{\mega\pascal}; deformation moduli in \si{\giga\pascal};
unit weight in \si{\kilo\newton\per\metre\cubed}; lengths in \si{\metre};
displacements in \si{\milli\metre}. Groundwater pressures are given in
\si{\kilo\pascal} or as an equivalent head in metres, with the datum stated.

\noindent\textbf{Dimensional inconsistency in empirical correlations.}\enspace
Several widely used correlations are not dimensionally homogeneous and are
therefore valid only in the units in which they were calibrated. Rock mass
modulus expressions returning \si{\giga\pascal} from a $\GSI$ value and a
$\sci$ in \si{\mega\pascal} are the common example. Wherever such an
expression appears in this book, the required input units are stated with it.
Substituting into these relations in unfamiliar units is a reliable way to
produce an answer wrong by orders of magnitude, and the error is not always
obvious in the result.

\noindent\textbf{Rounding and reported precision.}\enspace Design parameters
are reported to a precision consistent with their derivation. A rock mass
strength derived from a $\GSI$ estimated to the nearest five points does not
support three significant figures, and quoting it that way misrepresents the
evidence. This applies with particular force to numerical output, which is
returned to machine precision regardless of the quality of its input.
